
\documentclass[a4paper,12pt]{article} 
%%%% choose one of the following 4 lines to set colour and/or  date style, 
%\usepackage{pa} % Day, num Month year, and colour (use lp2)
\usepackage{fancyhdr}
\usepackage[a4paper, total={6in, 8in}]{geometry}
%\usepackage{fullpage}
\fancyhead{}
\fancyhead[RO,RE]{\bf Bruce A. Desmarais}
\fancyhead[LO,LE]{\bf Teaching Statement}


\begin{document} 
\pagestyle{fancy}

%\vspace*{-.8cm}

% Dissemination of information 
% Rigorous, objective, analytical standpoint
% Critical and analytical skill-building
%\vspace{1.25cm}


I have taught graduate and undergraduate level courses in research methods and data science throughout my career. I thrive on teaching in interdisciplinary quantitative and computational methods contexts, as the cross-disciplinary emphasis permits a sharp focus on the universal properties of the methods under study. I have taught courses in undergraduate and graduate introductory statistical methods, network analysis, machine learning, research design for social data analytics, a graduate survey of advanced methods in social data analytics, and a graduate course on the ``data layer'' in data analytics---emphasizing data collection, management, and cleaning. Instruction is the best mechanism by which our collective knowledge as scientists can be disseminated outside of our profession, or, in the case of graduate instruction, to those just entering the profession. I would welcome the opportunity to continue developing my instructional offerings relevant to computational social science, artificial intelligence, and democracy.

To tackle complex and technical data science subject-matter, I build courses based on regular and multimodal student activities. These include structured in-class discussion, regular problem sets, student-led code tutorials, and incrementally developed research projects. I use active learning segments in every course to provide students with strong incentives to study regularly throughout the semester, and to give them agency in crafting the course content. I find that constituting the classroom as an environment in which students learn by using the course material offers three concrete instructional advantages. First, using information requires students to develop a thorough degree of comprehension. Second, attempting to use information provides an intimate self-check on progress in learning. Third, because most active learning exercises are group-based, students are exposed to diverse ideas offered by their classmates and are motivated to maintain coverage of the material that is comparable to others in the class. One exercise that I use frequently in my research methods course is to co-develop a code tutorial with each student in the class, to be delivered by the student during class. This is helpful not only for the student developing the tutorial, but also builds confidence and collegiality among the other students in the class, since they see that one of their peers can master the methodology. 

My commitment to encouraging students to ``learn through doing'' manifests in graduate courses as doing everything I can to help students develop complete original research contributions in the context of the course. As is common, I require students to write an original research paper for each graduate course. I go beyond that by requiring students to complete their data analysis homework assignments using the data for their papers and supply me with the data and heavily commented code to demonstrate that they had correctly completed the assignments. Keeping the material closely connected to students' research objectives, and requiring repeated use and demonstration of mastery in-class, represent core components of my methodological instruction. It requires extra work to go through each student's individual application while grading assignments, but I find this to be a rewarding use of my time, and by the end of the semester I am deeply familiar with each student's project. 

My instructional perspective as a computational social scientist deserves special attention. I am constantly working to improve my performance as a quantitative methods instructor. Graduate quantitative methods instruction in political science and the social sciences more generally tends to present the widest variation in students' orientations towards the material---from students eager to understand the theoretical underpinnings of statistical tools to those who are shocked that graduate education in political science requires intensive statistics training. I find it to be particularly critical to take multi-pronged and multi-layered instructional approaches. To communicate any concept, I design lecture materials that are both intuitive and theoretically intensive---layered such that all students have access to the necessary concepts and the more technically inclined can satisfy their interests, present published applications that illustrate the use of the method, walk through examples in software and supervise implementation by students. I also provide numerous checkpoints at which students can evaluate their progress in mastering the material. 

The future of education related to artificial intelligence lies in its increasing relevance across disciplines. As a field, artificial intelligence is poised for growth like almost no other. The challenge for programs in this space going forward will be to develop strategies to embrace and capitalize on the breadth and diversity of the potential future student body. My approach to data science instruction reflects the movement for, ``Data Science for Everyone.'' This is a philosophy of data science instruction that establishes a top priority of making the coursework accessible to all who are interested in learning, regardless of their math or computing background, and with careful attention to how individuals' identities may affect their experience of or perspective on the course content. With each new semester I revise my syllabus in ways that lower barriers to entry, while maintaining high standards of rigor. One of the best ways that we can advance the ethics of artificial intelligence development and application is foster diversity, equity, and inclusion in training programs. 

I am passionate about and dedicated to building institutions that facilitate sustained interdisciplinary engagement in the field of data science. At UMass Amherst, I worked with colleagues to transform an initiative in CSS into robust institutions that would develop into the Computational Social Science Institute and a masters degree program in Data Analytics and Computational Social Science. When I arrived at Penn State in 2015, there was NSF funding to support the development of graduate training in Big Data Social Science. I have taken a leading role in developing this investment into a Center for Social Data Analytics and a dual-title PhD program in Social Data Analytics, both of which I now direct. I would be thrilled to contribute to the development and maintenance of interdisciplinary institutions related to artificial intelligence and computational social science at McGill. 





\end{document} 
